\subsection{D-brane Instantons in Type II String Theory --- arXiv:0902.3251}

\textbf{Authors:} Ralph Blumenhagen, Mirjam Cvetic, Shamit Kachru, Timo Weigand

\paragraph{主な主張}
四次元$\mathcal{N}=1$を保つtype IIコンパクト化で、D-brane instantonの零モードがsuperpotentialなどの正則結合への寄与を決めることを整理する~\cite{Blumenhagen:2009qh}。

\paragraph{新規性}
多instanton・背景フラックス・特異点上のquiverを含む計算法と、物質結合や超対称性破れへの応用を統一的にまとめる。

\paragraph{位置付け}
モジュライ安定化や摂動論で禁じられた物質結合の生成を扱うためのD-brane instantonの総説。

\paragraph{コメント（読書メモ）}
本人の読書メモ（2024年4月20日）：ゲージ理論のinstantonはEuclidean自己双対方程式$F=*F$の解として現れる。弦理論については、完全な第二量子化の定式化が未完成だという本レビュー執筆時の説明を踏まえ、場の理論との類推を非結合極限や双対性で検証する、という理解を記した。

ヘテロティック弦で早くから研究されたworldsheet instantonは、内部多様体の非自明な二次元サイクルを包むEuclidean閉弦であり、四次元では局在した対象になる。元メモの「2階巻きつく」は二次元サイクルを包む意味として整理する。その効果は弦結合$g_s$ではなく世界面の展開パラメータ$\alpha'$に関して非摂動的であり、$\alpha'$の名称については当時「ledge slope ?」と疑問を残していた。

本レビューが述べる約15年間の進展として、時空の観点からも非摂動的なD-braneが注目され、四次元時空を満たし内部サイクルを包む交差D-brane上のゲージ場・物質スペクトルから、現実的模型と$\mathcal{N}=1$四次元有効超重力の構成が進んだと理解した。

さらに、D-braneは模型構築の背景となるだけでなく非摂動効果にも寄与する。内部の非自明なサイクルだけを包むEuclidean D$p$-braneは四次元時空の一点に局在するため、D-brane instantonまたはE$p$-braneと呼ばれる。

その指数抑制は一般に四次元のゲージ結合だけでは決まらず、stringy/exotic instantonと呼ばれる場合がある、という点を記録した。元メモでは四次元ゲージ結合を$g_s$と記していたが、ここでは弦結合$g_s$との混同を避けて区別する。
